\section*{Chapter \ref{ch:b3:complex-kinetics} --- Complex Kinetics: Chains, Enzymes and Oscillations}

\begin{solution}{exo:b3:complex-kinetics:1}
$\ce{Cl2 -> 2 Cl}$: initiation; $\ce{Cl + H2 -> HCl + H}$ and $\ce{H + Cl2 -> HCl + Cl}$:
propagation; $\ce{2 Cl -> Cl2}$: termination. Carriers: Cl and H. Overall (sum of the
two propagation steps): $\ce{H2 + Cl2 -> 2 HCl}$.
\end{solution}

\begin{solution}{exo:b3:complex-kinetics:2}
$[\mathrm M]_{1/2} = k_2/k_{-1} = \qty{1.0e-3}{mol/L} = \qty{1.0}{mol/m^3}$; $p = cRT = 1.0 \times 8.314 \times 750 =
\qty{6.2}{kPa}$.
\end{solution}

\begin{solution}{exo:b3:complex-kinetics:3}
$v/V_{\max} = [\mathrm S]/(K_M + [\mathrm S])$: $\frac12$, $\frac34$, $\frac{9}{10}$. 90\,\% of $V_{\max}$ needs $[\mathrm S] = 9K_M$.
\end{solution}

\begin{solution}{exo:b3:complex-kinetics:4}
$k_{\mathrm{cat}}/K_M = 100/\num{0.80e-3} = \qty{1.25e5}{L.mol^{-1}.s^{-1}}$: some 60\,000 times below the diffusion
limit, and close to the typical value.
\end{solution}

\begin{solution}{exo:b3:complex-kinetics:5}
Steady state for \ce{CH3CO}: $k_2[\ce{CH3}][\ce{CH3CHO}] = k_3[\ce{CH3CO}]$. Steady state for \ce{CH3}:
$k_1[\ce{CH3CHO}] - k_2[\ce{CH3}][\ce{CH3CHO}] + k_3[\ce{CH3CO}] - 2k_4[\ce{CH3}]^2 = 0$. Adding: $k_1[\ce{CH3CHO}] =
2k_4[\ce{CH3}]^2$, so $[\ce{CH3}] = (k_1/2k_4)^{1/2}[\ce{CH3CHO}]^{1/2}$, and $\dd[\ce{CH4}]/\dd t = k_2[\ce{CH3}][\ce{CH3CHO}] =
k_2(k_1/2k_4)^{1/2}[\ce{CH3CHO}]^{3/2}$.
\end{solution}

\begin{solution}{exo:b3:complex-kinetics:6}
At half conversion $[\ce{H2}] = [\ce{Br2}] = \frac12$, $[\ce{HBr}] = 1$ (relative units):
$v/v_0 = \frac12 \times (\frac12)^{1/2}/(1 + 0.10 \times 1/0.5) = 0.354/1.2 = 0.29$. Without inhibition it would
be 0.354: inhibition divides the rate by a further 1.2.
\end{solution}

\begin{solution}{exo:b3:complex-kinetics:7}
Explosion when $1000p > 6000/p + 10p^2$, i.e. $10p^3 - 1000p^2 + 6000 < 0$, whose positive
roots are 2.48 and \qty{99.9}{kPa}: the mixture explodes between about 2.5 and
\qty{100}{kPa} (first and second limits of the model).
\end{solution}

\begin{solution}{exo:b3:complex-kinetics:8}
$1/v = 1/V_{\max} + (K_M/V_{\max})/[\mathrm S]$: $5.0 = 1/V_{\max} + 2.0K_M/V_{\max}$ and $2.5 = 1/V_{\max} + 0.5K_M/V_{\max}$.
Subtracting: $K_M/V_{\max} = \qty{1.667}{s.mM.\micro M^{-1}}$, then $1/V_{\max} = 1.667$: $V_{\max} = \qty{0.60}{\micro M.s^{-1}}$,
$K_M = \qty{1.0}{mM}$.
\end{solution}

\begin{solution}{exo:b3:complex-kinetics:9}
Same intercept on the $1/v$ axis: competitive. $K_M^{\mathrm{app}} = 1/0.417 = \qty{2.40}{mM} = 3K_M$, so
$1 + 50/K_i = 3$ and $K_i = \qty{25}{\micro M}$.
\end{solution}

\begin{solution}{exo:b3:complex-kinetics:10}
$N = \qty{1.01}{mol/L}$: $t_{1/2} = \ln(1.01/0.010 - 1)/(0.50 \times 1.01) = \ln100/0.505 = \qty{9.1}{s}$. The
rate $kx(N - x)$ is largest at $x = N/2$: the same instant, the inflexion of the S
curve.
\end{solution}

\begin{solution}{exo:b3:complex-kinetics:11}
Trace $B - 2$, determinant 1. $B = 1.5$: $\lambda = -0.25 \pm 0.968\iu$, a stable focus (damped
oscillations towards the steady state). $B = 2.5$: $\lambda = 0.25 \pm 0.968\iu$, an unstable focus:
oscillations grow until they reach the \omterm{def:b3:complex-kinetics:oscillating}{limit cycle}.
\end{solution}

\begin{solution}{exo:b3:complex-kinetics:12}
$1/\tau = \num{1.0e8} \times \num{5.40e-5} + \num{1.0e3} = \qty{6.4e3}{s^{-1}}$, $\tau = \qty{156}{\micro s}$. Measure $\tau$ at several
total concentrations: $1/\tau$ against $[\mathrm A]_e + [\mathrm B]_e$ is a straight line of slope $k_1$ and
intercept $k_{-1}$.
\end{solution}

\begin{solution}{pb:b3:complex-kinetics:1}
\textbf{1.} At the start $[\mathrm S]$ is known and the product, which could inhibit or react
back, is absent.
\textbf{2.} The low-$[\mathrm S]$ points, at large $1/[\mathrm S]$ and $1/v$: they set the slope, and the
inversion magnifies their errors.
\textbf{3.} It weights each measured rate as measured; the Lineweaver--Burk line gives
biased parameters.
\textbf{4.} $K_M$: $0.773$ is 1.2 standard errors below 0.801; $V_{\max}$: $0.491$ is 2.2 standard
errors below 0.502. Both are low, as expected from the bias.
\textbf{5.} $k_{\mathrm{cat}} = 0.502/0.0050 = \qty{100}{s^{-1}}$.
\textbf{6.} $100/\num{0.801e-3} = \qty{1.25e5}{L.mol^{-1}.s^{-1}}$.
\textbf{7.} Some 60\,000 times below the diffusion limit; a typical \omterm{def:b3:complex-kinetics:enzyme}{enzyme}.
\textbf{8.} Competitive: $V_{\max}$ unchanged, apparent $K_M$ growing with $[\mathrm I]$.
\textbf{9.} $K_M^{\mathrm{app}} = K_M(1 + [\mathrm I]/K_i) = K_M + (K_M/K_i)[\mathrm I]$.
\textbf{10.} $K_i = 0.799/0.0321 = \qty{24.9}{\micro M}$.
\textbf{11.} $24.9 \pm 4.30 \times 0.9 = 24.9 \pm \qty{3.9}{\micro M}$.
\textbf{12.} The dissociation constant of the enzyme--inhibitor complex: the
inhibitor concentration that occupies half the free \omterm{def:b3:complex-kinetics:enzyme}{enzyme}.
\textbf{13.} $v_i/v_0 = (K_M + [\mathrm S])/(K_M(1 + [\mathrm I]/K_i) + [\mathrm S])$.
\textbf{14.} $2/3 = 0.67$.
\textbf{15.} $2/12 = 0.17$.
\textbf{16.} $11/21 = 0.52$: at high substrate concentration the substrate outcompetes
the inhibitor.
\textbf{17.} At $[\mathrm S] = K_M$, $v_i/v_0 = 2/(2 + [\mathrm I]/K_i) = 0.10$ gives $[\mathrm I]/K_i = 18$.
\textbf{18.} $18 \times 24.9 = \qty{448}{\micro M}$.
\textbf{19.} The first step is fast: every \omterm{def:b3:complex-kinetics:enzyme}{enzyme} molecule quickly releases one $\mathrm P_1$
and is trapped as the \omterm{def:b3:complex-kinetics:enzyme}{acyl-enzyme}; afterwards $\mathrm P_1$ appears only as fast as the slow
second step frees the \omterm{def:b3:complex-kinetics:enzyme}{enzyme}.
\textbf{20.} $1.28/2.0 = 0.64 = (k_2/(k_2 + k_3))^2$, so $k_2/(k_2 + k_3) = 0.80$.
\textbf{21.} $200/2.0 = \qty{100}{s^{-1}}$, the same as in Part I.
\textbf{22.} $k_2 = 0.80 \times 625 = \qty{500}{s^{-1}}$, $k_3 = \qty{125}{s^{-1}}$.
\textbf{23.} $500 \times 125/625 = \qty{100}{s^{-1}}$.
\textbf{24.} \textbf{$[\mathrm I]_{90\,\%} = 18K_i \approx \qty{0.45}{mM}$ at $[\mathrm S] = K_M$.}
\end{solution}
